\documentclass[10pt]{article}
\usepackage[margin=20mm]{geometry}
\usepackage{amsmath,amssymb,mathtools,fontspec,unicode-math,xurl,hyperref}
\IfFontExistsTF{DejaVu Serif}{\setmainfont{DejaVu Serif}}{\setmainfont{Latin Modern Roman}}
\setmathfont{latinmodern-math.otf}
\setlength{\parindent}{0pt}\setlength{\parskip}{6pt}\sloppy
\title{Initial scientific review of
Persistence and ball exponents for Gaussian stationary processes}
\author{AI-assisted technical review}
\date{}
\begin{document}
\maketitle
Abstract: This review concerns the acquired arXiv version 2112.04820v3 (3 April 2025), by Naomi D. Feldheim, Ohad N. Feldheim and Sumit Mukherjee, associated with DOI 10.1002/cpa.22255. Six major issues concern false auxiliary estimates or invalid derivations used in the principal arguments, including the advertised example of nonexistence. A further minor issue concerns the stated ranges of applications and remarks. Explicit counterexamples and local repair arguments are given below. No counterexample to Theorems 1–5 is established here. The scientific decision remains pending; completion of this initial review does not constitute independent acceptance.\par
\section*{Version, scope and mathematical framework}
All 51 pages of the acquired source, including its references, were read in this conversation. Page numbers below refer to that PDF, and equation numbers prefixed by source refer to the paper. The journal version of record was identified but was not acquired for a line comparison. The current arXiv record states that v3 corrected earlier inaccuracies in Lemmas 1.1 and 1.3. The issues below were checked in v3 itself. The publisher records a separate correction dated 26 June 2025 replacing the infimum by the supremum in the abstract’s ball-exponent definition. That acknowledged correction is not counted as a new finding.\par
Equation reference: eq:def. The covariance convention used here is \[r_\rho(t)=\int_{\mathbb R}e^{-it\lambda}\,d\rho(\lambda),\qquad
 \theta_\rho^\ell(T)=-T^{-1}\log\mathbb P(\inf_{[0,T]}f_\rho>\ell).
 \] The spectral measure is finite, nonnegative and symmetric. The class \(\mathcal L\) imposes a finite logarithmic moment of order \(1+\beta\). The paper explicitly defines \(\mathcal M\) by \(\rho'(0)=\lim_{\varepsilon\downarrow0}
\rho([-\varepsilon,\varepsilon])/(2\varepsilon)\in(0,\infty]\). Thus an infinite density, including an atom at zero, is allowed. Uniform assertions use fixed \(\alpha,\beta,B,\ell\) and local lower or upper bounds on these averaged densities. The finite positive density case, the infinite density case, and the zero density case must be kept distinct. Ball probabilities use \(\sup|f_\rho|\) and radius \(\ell>0\).\par
Theorem 1 asserts existence at every real level when \(\rho\in\mathcal L\cap\mathcal M\), and, with a nonzero absolutely continuous component, positivity precisely when \(\rho'(0)<\infty\). Theorem 2 concerns the vanishing ball exponent of singular measures. Theorems 3–5 concern spectral perturbations, continuity and sampling. The main persistence argument uses the finite-time comparison lemmas, the singular ball result, and smooth approximation in a specific order. This dependency structure determines the severity assignments below; repairability does not make an invalid essential derivation valid.\par
\section*{Global level bound and the Gaussian change of measure}
Problem 1 (P01: Major).  Lemma 1.1, p. 7, claims a bound linear in every \(\delta>0\) with a constant independent of \(\delta\). The change-of-measure calculation on p. 19 omits a necessary quadratic cost and uses the likelihood ratio in the wrong direction.\par
\subsection*{Hypotheses and decisive counterexample}
Equation reference: eq:quadratic. Take \(d\mu=(2\pi)^{-1}\mathbf 1_{[-\pi,\pi]}d\lambda\), \(\ell=0\), \(T=4\). This compactly supported measure belongs to \(\mathcal L_{\beta,B}\), for example for \(\beta=1\) and sufficiently large fixed \(B\), and to \(\mathcal M_{\alpha,1}\) for \(\alpha=1/(4\pi)\). Its covariance is \(\sin(\pi t)/(\pi t)\), so \(f_\mu(0),\ldots,f_\mu(4)\) are independent standard normal variables. Writing \(Q(x)=\mathbb P(Z>x)\), \[\theta_\mu^\delta(4)\ge-\frac54\log Q(\delta)
       =\frac58\delta^2+O(\log\delta).
 \] By continuity of sample paths on the fixed interval, \(\mathbb P(\inf_{[0,4]}f_\mu>-\delta)\to1\), whence \(\theta_\mu^{-\delta}(4)\to0\). Consequently the difference in Lemma 1.1 cannot be at most \(C\delta+\log(2)/4\) for all \(\delta\).\par
An independent evaluation of the corresponding discrete five-point probabilities gives differences \(43.76680\), \(164.61925\), and \(645.48206\) at \(\delta=8,16,32\). These are exact Gaussian-tail formula evaluations, not measurements of continuous-path probabilities. The analytic argument [eq:quadratic] (\#eq:quadratic) establishes the continuous counterexample without that distinction being blurred.\par
\subsection*{Proof defect, alternative reading and repair}
Equation reference: eq:rn, eq:locallevel. On pp. 18–19 the coefficient \(\zeta\) is shifted to mean \(-a\), where \(a=c\delta\sqrt T\). If \(p\) and \(q\) are the respective densities of \(N(0,1)\) and \(N(-a,1)\), the required ratio is \[\frac{q(x)}{p(x)}=\exp(-ax-a^2/2).
 \] Restricting the shifted coefficient’s integration variable to \(|x|\le\sqrt{MT}\), and using the paper’s tail truncation, therefore gives \[\theta_\mu^{\ell+\delta}(T)-\theta_\mu^{\ell-\delta}(T)
 \le c\sqrt M\,\delta+\frac{c^2}{2}\delta^2+\frac{\log2}{T}.
 \] Here \(M\) can be fixed from the lower persistence bound at level \(\ell\), because the lower event at level \(\ell-\delta\) contains that event. The printed \(p/q\) estimate cannot remove the term \(a^2/2\).\par
For \(0<\delta\le\delta_0\), [eq:locallevel] (\#eq:locallevel) is linear with a constant depending on \(\delta_0\). This supplies the relevant local form once the lower probability input is justified as discussed in P02. The uses in Lemma 1.3 and Theorem 4 involve small level changes. Thus the example refutes the global auxiliary assertion, but does not refute local Lipschitz continuity of the limiting exponent. The proof of Theorem 4 on p. 34 explicitly depends on Lemma 1.1; the same input recurs in the proofs of Theorems 1, 3 and 5. That essential role, rather than a claim of failure of all final results, is the reason for the Major classification.\par
\section*{An increment estimate cannot bound the absolute supremum}
Problem 2 (P02: Major).  Lemma 2.19, pp. 14–15, bounds \(\mathbb E\sup_{[0,h]}|f_\rho|\) by a positive power of \(\sup_{[0,h]}|r_\rho(0)-r_\rho(t)|\). The proof of Lemma 2.22, pp. 16–17, then makes this absolute supremum arbitrarily small by shortening the interval. Both steps fail.\par
\subsection*{Counterexample and dependency}
For \(\rho=\delta_0\), which satisfies the logarithmic moment assumption, \(f_\rho(t)=Z\) at every time. The proposed right-hand side is zero, whereas \[\mathbb E\sup_{[0,h]}|f_\rho|=\sqrt{2/\pi}>0.\] The issue is not limited to degenerate processes: for every continuous process of nonzero variance, the left side is at least \(\mathbb E|f_\rho(0)|>0\), while the proposed right side tends to zero as \(h\downarrow0\).\par
For the constant process and radius \(0.1\) the ball probability on every nonempty interval is \(\mathbb P(|Z|<0.1)=0.07965567\). Thus the intermediate conclusion in Lemma 2.22 that a sufficiently short interval has ball probability at least \(1/2\) is also false. Dudley’s bound stated in Proposition 2.15 controls a centered supremum or anchored increments; adding absolute values requires accounting for the base value. The absolute values are present in the actual source images, including the downstream proof.\par
Lemma 2.22 supplies Lemma 2.23, the uniform lower persistence bound. That bound enters the likelihood truncation in P01, Proposition 3.1, Lemma 5.7 and the comparison arguments for Theorems 1 and 4. The false estimate therefore invalidates an operative proof chain.\par
\subsection*{A local correction preserving the claimed lower bound}
Equation reference: eq:smallballrepair. Set \(Y_t=f_\rho(t)-f_\rho(0)\). The logarithmic increment estimate and Dudley’s entropy integral give, with a constant depending on \(\beta,B\), \[\mathbb E\sup_{[0,h]}|Y_t|\le C_{\beta,B}
       \bigl(\log(1/h)\bigr)^{-\beta/2}.\] This is an anchored estimate; it does not assert that \(\mathbb E|f_\rho(0)|\) vanishes. Choose \(h<1\) so that its right side is at most \(\ell/4\). Markov’s inequality gives \(\mathbb P(\sup|Y|<\ell/2)\ge1/2\). The Gaussian correlation inequality [royen], applied first to finite collections of values and then by continuity, gives \[\begin{aligned}
 \mathbb P(\sup_{[0,h]}|f_\rho|<\ell)
 &\ge \mathbb P(|f_\rho(0)|<\ell/2,\ \sup_{[0,h]}|Y|<\ell/2)\\
 &\ge \tfrac12\mathbb P(|N(0,2B)|<\ell/2)=:q_{\ell,B}>0 .
 \end{aligned}\] Indeed \(\rho(\mathbb R)\le2B\), and both events are symmetric convex constraints on the same centered Gaussian vector; no independence is assumed. Repeated Gaussian correlation and stationarity give \[\mathbb P(\sup_{[0,T]}|f_\rho|<\ell)\ge q_{\ell,B}^{\lceil T/h\rceil}
       \ge e^{-C_{\ell,\beta,B}T},\qquad T\ge1 .\] This proves the needed conclusion of Lemma 2.22 for \(\ell>0\) with a different short-interval constant. Likewise the estimate needed for Lemma 2.20 can retain \(\mathbb E|f_\nu(0)|\le\sqrt{\nu(\mathbb R)}\); bounded total mass allows this term to be absorbed into the stated fractional power of \(\nu(\mathbb R)\). These repairs are counterevidence to any assertion that the final small-ball or persistence theorems have been disproved.\par
\section*{The uniform threshold in the concentration lemma}
Problem 3 (P03: Major).  Lemma 2.21, pp. 15–16, asserts that variance at most \(\varepsilon/T\) and increment variance at most \(a^2h^2\) imply \(\mathbb P(\sup_{[0,T]}|W|>\delta)\le2e^{-\delta^2T/(8\varepsilon)}\) whenever \(T>c(a^{-1}+\delta^{-1}+1)\), with a universal \(c\). The stated threshold is false.\par
\subsection*{A smooth stationary counterexample}
Equation reference: eq:entropycontradiction. Fix \(\varepsilon=1/2\), \(a=1\), and \(K=32\). For each \(\delta>0\), put \[T=16\delta^{-2},\qquad \sigma^2=\varepsilon/T=\delta^2/32,\qquad
 D=\sqrt{96}/\delta .\] Let \(W\) have the compactly supported spectral density \(\sigma^2(2D)^{-1}\mathbf 1_{[-D,D]}\). Then \(\operatorname{Var}W=\sigma^2\) and, using \(2(1-\cos x)\le x^2\), \[\operatorname{Var}(W(t+h)-W(t))\le \frac{\sigma^2D^2}{3}h^2=h^2.\] At the points \(j\pi/D\), the values are independent \(N(0,\sigma^2)\). There are \(N=\lfloor TD/\pi\rfloor+1\) such points in \([0,T]\). Consequently, with \(p=\mathbb P(|Z|>\sqrt{32})>0\), \[\mathbb P(\sup_{[0,T]}|W|>\delta)\ge1-(1-p)^N\longrightarrow1
       \quad(\delta\downarrow0).
 \] The claimed upper bound is the constant \(2e^{-4}=0.03663128\). Also \(T/(a^{-1}+\delta^{-1}+1)=16/(\delta+2\delta^2)\to\infty\), so the hypothesis holds for any proposed universal \(c\) when \(\delta\) is sufficiently small. This addresses the unknown constant, not merely one selected numerical value of it.\par
The independently evaluated lower bounds are \(0.53667678\) at \(\delta=0.01\) and greater than \(0.999\) at \(\delta=0.001\). The computation uses the exact independent-grid formula; it neither simulates paths nor allocates a vector of \(N\) samples.\par
\subsection*{Cause, downstream uses and repair boundary}
The inference to source (2.5) loses the necessary metric-entropy scale. A valid sufficient estimate, with universal constants, is \[\mathbb E\sup_{[0,T]}|W|\le
 C\sigma\sqrt{\log(2+aT/\sigma)},\qquad \sigma^2=\varepsilon/T.\] This follows by retaining a base point and integrating the covering bound up to diameter at most \(2\sigma\). To use Borell with a deviation \(\delta/2\), it suffices to require \(\delta\ge2C\sigma\sqrt{\log(2+aT/\sigma)}\).\par
The incorrect lemma is explicitly used on p. 22 in Proposition 3.1 and on pp. 39–41 in Proposition 5.3. Those results feed spectral continuity and existence in Theorems 1 and 4. Their particular processes have additional structure: on p. 22 the actual bounds are \(\sigma^2\le2/(L^2T)\) and \(a_T^2\le2/T^3\). The corrected entropy estimate then gives \[\mathbb E\sup|g|\le
 C\frac{\sqrt{\log(2+L)}}{L\sqrt T},\] which is compatible with the chosen large \(L\) and \(T\). Proposition 5.3 also states a stronger threshold involving \(\varepsilon^{-2}\). These facts make repair plausible, but all its uniform constants and parameter dependencies must be checked. The general lemma itself cannot be retained as stated, and this review does not certify a completed repair of every downstream estimate.\par
\section*{The truncation argument removes a non-small amount of mass}
Problem 4 (P04: Major).  In Proposition 3.1, pp. 21–23, the passage after source (3.9) uses \((1-h^2)\nu_{L/T}(\mathbb R)\le1/T\), claiming this follows when \(L\ge1/(2\alpha')\). The inequality is false even within the proposition’s full parameter range, and the invocation of Lemma 1.3(I) therefore lacks its small-mass hypothesis.\par
\subsection*{All parameter restrictions and a counterexample}
The proposition assumes \(\mu\in\mathcal M_{\alpha,A}\cap\mathcal L_{\beta,B}\), \(\nu\in\mathcal M_{(0,\alpha'),D}\cap\mathcal L_{\beta,B}\), \(0<D\le1\), \(0<\alpha'\le B\), \(0<\eta<c_1\alpha'\), \(T\ge\max(c_2/\eta,1/A)\) and \(c_3/\eta\le L\le\alpha'D^2T/B\). Choose \(\mu\) as in P01, \(A=\beta=1\), \(B=10\), \(d\nu=\mathbf 1_{[-1/2,1/2]}d\lambda\), \(\alpha'=2\), and \(D=1/2\). Set \(L=\eta^{-2}\) and \(T=160L\). For every fixed choice of the proposition’s constants all its assumptions hold when \(\eta\) is sufficiently small. In particular \(\alpha'D^2T/B=8L\).\par
Equation reference: eq:removedmass. With \(M=L^{1/4}\) and the paper’s normalized \(\operatorname{sinc}u=\sin(\pi u)/(\pi u)\), its actual removed mass is \[\begin{aligned}
 \Delta_L
 &=2\int_0^{L/T}\{1-\operatorname{sinc}^4(T\lambda/M)\}\,d\lambda\\
 &=\frac{2}{160}
       -\frac{2M}{T}\int_0^{L/M}\operatorname{sinc}^4 u\,du
       \longrightarrow\frac1{80}.
 \end{aligned}\] For a completely elementary bound, \(\int_0^\infty\operatorname{sinc}^4u\,du\le1+1/(3\pi^4)\). Thus the limit and a positive lower bound follow without numerical quadrature. In contrast \(1/T\to0\). At \(L=10^4\), evaluation of this rigorous lower-bound expression gives \(0.01248746\), whereas \(1/T=6.25\times10^{-7}\).\par
The intermediate bound \(\nu([-L/T,L/T])\le2\alpha'L/T\) is valid. It implies \(2\alpha'L/T\le1/T\) only for \(L\le1/(2\alpha')\), the opposite inequality. Merely reversing that condition would conflict with the required large-\(L\) regime. Equation [eq:removedmass] (\#eq:removedmass) also rules out interpreting the mistake as an innocuous loose estimate of a genuinely tiny removed mass.\par
\subsection*{Dependency and a way to avoid the false step}
Proposition 3.1 is the truncation input to Lemma 1.3(II), then Proposition 5.2 and the proofs of Theorems 1, 3 and 4. Its conclusion has not been disproved.\par
A local modification is to apply Lemma 1.2 with base measure \(\mu+\nu_{L/T}\) and smoothed measure \(\nu-\nu_{L/T}\). This gives directly, with the same shortened interval \(I_T\), \[\mathbb P(f_{\mu+\nu}>\ell\text{ on }[0,T])
 \le
 \mathbb P(f_{\mu+\nu_{L/T}+h^2(\nu-\nu_{L/T})}>\ell
       \text{ on }I_T).\] The high-frequency remainder is exactly the one estimated on p. 22, while the low-frequency measure already has the desired form. No subsequent removal of \((1-h^2)\nu_{L/T}\) is necessary. The base retains the same lower local-density bound and has a logarithmic-moment bound with \(2B\) in place of \(B\). This modification removes the demonstrated obstruction. Completion still requires the corrected concentration and local-level bounds in P03 and P01 with appropriate uniform constants.\par
\section*{The spectral projection fails for an atom at zero}
Problem 5 (P05: Major).  Lemma 4.2, pp. 25–26, allows every spectral measure supported on \([-D,D]\) but assigns zero asymmetrically to the first positive bin. The independent decomposition in source (4.2) is false when that bin contains both a zero atom and nonzero spectral mass.\par
\subsection*{Exact covariance contradiction}
Take \(D=n=1\) and \[\rho=\tfrac13(\delta_{-1}+\delta_0+\delta_1).\] This is a finite symmetric nonnegative measure, including a legitimate purely singular measure for Proposition 4.1. The source defines \(I_1=[0,1]\), \(I_{-1}=[-1,0)\), hence \[C_1(t)=\frac{1+\cos t}{2},\qquad S_1(t)=\frac{\sin t}{2},
 \qquad \rho(I_1\cup I_{-1})=1.\] Its asserted representation would be \(f_\rho(t)=\zeta C_1(t)+\eta S_1(t)+R(t)\), with \(R\) independent of the standard normal coefficients. At \(t=0\) the first term already has variance one, equal to \(\operatorname{Var}f_\rho(0)\). Therefore \(\operatorname{Var}R(0)=0\). At \(t=\pi\) both \(C_1\) and \(S_1\) vanish. Independence consequently forces \(\operatorname{Cov}(f_\rho(0),f_\rho(\pi))=0\). The true spectral covariance is \[r_\rho(\pi)=\tfrac13(1+2\cos\pi)=-\tfrac13,\] a contradiction.\par
The proof on p. 26 treats the positive and negative bins as equal-mass partners in its Hilbert-space normalization. Here their masses are \(2/3\) and \(1/3\). This cannot be fixed by assigning a value to a zero-denominator expression: both relevant denominators are nonzero.\par
\subsection*{Repair and impact boundary}
Separate \(\sqrt{\rho(\{0\})}\,Z_0\) as its own independent constant component. Partition the remaining strictly positive spectrum and its reflections into equal-mass pairs, omitting empty bins. The ordinary orthogonal projection then gives the cosine and sine terms and residual estimates for the remaining measure. The additional constant Gaussian component has rank one. It does not affect the sublinear-rank argument in Proposition 4.1: its contribution to a fixed-radius ball event has a positive \(T\)-independent probability. The error is therefore repairable.\par
As printed, however, Lemma 4.2 is the explicit input at p. 27 to Proposition 4.1, which supplies Theorem 2. The singular ball theorem also enters the persistence approximation on p. 31 and the invariance statements in Theorem 3. Measures with an atom at zero are not excluded there. This invalid core decomposition is Major; it is not a counterexample to the corrected singular ball theorem. The applications in Proposition 5.3 have additional local mass restrictions, so the counterexample is not asserted to violate that proposition’s hypotheses.\par
\section*{The nonexistence example has a finite positive origin density}
Problem 6 (P06: Major).  Remark 1, p. 2, and Section 5.4, pp. 41–42, advertise nonexistence of the zero-level persistence exponent for \[w_{a,b}(\lambda)=
 \left(\frac{a+b}{2}+\frac{b-a}{2}\cos(1/\lambda)\right)
 \mathbf 1_{\{|\lambda|\le1\}},\qquad 0<a<b.\] Source (5.19) claims that the averaged density has lower and upper limits \(a\) and \(b\). In fact its limit is \((a+b)/2\).\par
Claim 1.  For every \(0<a<b\), the displayed measure belongs to \(\mathcal L\cap\mathcal M\) and has \(\rho'(0)=(a+b)/2\).\par
Equation reference: eq:correctdensity. Proof. With \(x>0\), \[I(x):=\int_0^x\cos(1/t)\,dt
      =\int_{1/x}^{\infty}\frac{\cos u}{u^2}\,du.\] Integration by parts gives \[I(x)=-x^2\sin(1/x)
       +2\int_{1/x}^{\infty}\frac{\sin u}{u^3}\,du,
 \qquad |I(x)|\le2x^2.\] It follows that \[\frac{\rho([0,x])}{x}
   =\frac{a+b}{2}+\frac{b-a}{2}\frac{I(x)}x
   =\frac{a+b}{2}+O((b-a)x).
 \] Symmetry gives the same limit in the definition of \(\rho'(0)\). The density is bounded, nonnegative and compactly supported, so the logarithmic moment condition also holds. ◻\par
For \(a=1,b=3\), independent integration-by-parts evaluations at \(x=0.01\) and \(0.001\) give averaged densities approximately \(2.00523288\) and \(1.99917425\). The saved truncation-error formulas are analytic errors for the finite asymptotic expansion; the floating evaluations are not certified interval-arithmetic enclosures. The proof [eq:correctdensity] (\#eq:correctdensity), rather than those decimals, is decisive.\par
\subsection*{What fails and what could replace it}
The low-density subsequence required in the final argument on p. 42 does not follow. For the selected range \(a<b/3\), the actual limiting average is greater than \(3a/2\), so the required eventual values at most \(3a/2\) are absent. The paper’s own Theorem 1 would give existence for this very measure. Thus the advertised example and that theorem cannot both hold as written. This review establishes the false average and invalid nonexistence derivation; it does not use the contested proof of Theorem 1 as an independent proof of the final exponent.\par
A possible replacement uses a slowly oscillating density, for example near zero \[\widetilde w(t)=\frac{a+b}{2}
    +\frac{b-a}{2}\cos\!\left(\log\log(e/|t|)\right).\] It has the desired averaged lower and upper limits. Indeed, writing \(L_x=\log(e/x)\) and \(t=xu\), \[\int_0^1
 \left|\log\frac{\log(e/(xu))}{\log(e/x)}\right|du
 \le\frac1{L_x}\int_0^1\log(1/u)\,du=\frac1{L_x}.\] The averaged cosine therefore differs from its value at \(x\) by \(O(1/L_x)\). This proves the required oscillation of the averages for a different density. It is not a correction of \(\cos(1/\lambda)\) by reinterpretation. The probability estimates and constants must then be reapplied to the replacement, using the actual assumptions of [ffn]. The sharpness illustration is a substantive stated contribution, and its supplied proof is invalid; this explains the Major classification.\par
\section*{Missing range restrictions in applications and remarks}
Problem 7 (P07: Minor).  Example 3 on p. 8 extends the zero-density conclusion to every real level. Remark 3 on p. 3 omits the finite-origin-density restriction when discussing positivity. Explicit elementary cases contradict those ranges, while the main theorem statements retain relevant distinctions.\par
\subsection*{The moving-average example at a negative level}
Set \(U_0=1\), \(U_1=-1\), and all other \(U_k=0\). The sequence is square summable, \(\widehat U\) is continuous near zero, and \(\sum U_k=0\). The process is \(X_j=Z_j-Z_{j-1}\) for independent standard normal \(Z_j\). At level \(\ell=-1\), if all \(N+2\) variables indexed from \(-1\) to \(N\) lie in \((-1/2,1/2)\), then \(X_j>-1\) for \(0\le j\le N\). Hence \[\mathbb P(X_j>-1,\ 0\le j\le N)\ge p^{N+2},\qquad
 p=\mathbb P(|Z|<1/2)=0.38292492.\] Therefore \[\limsup_{N\to\infty}\theta_X^{-1}(N)
       \le-\log p=0.95991634<\infty.\] This disproves the displayed equivalence \(\theta_X^\ell=\infty\Longleftrightarrow\sum U_k=0\) at the stated all-real-level scope. No assumption about the existence of the limit is needed to refute an asserted infinite limit. Remark 2 and Lemma 2.24 correctly restrict the zero-density infinite-exponent conclusion to \(\ell\ge0\). Restoring that restriction in the application removes this contradiction.\par
\subsection*{Infinite origin density in Remark 3}
The paper’s \(\mathcal M\) includes \(\rho'(0)=\infty\). Already \(\rho=\delta_0\) has \(\theta_\rho^\ell=0\) for every finite \(\ell\), since its persistence probability is the constant \(Q(\ell)\). Thus the positive-level conjecture in that remark cannot include the whole class as literally defined.\par
The nonpositive-level assertion that positivity is equivalent to a nonzero absolutely continuous part also fails at infinite origin density. Let \(f=Z+g\), with \(Z\) independent standard normal and \(g\) an independent unit-variance bandlimited stationary process with nonzero absolutely continuous spectrum. Put \(A_0=\mathbb E\sup_{[0,1]}|g|<\infty\). A union of unit-interval Borell bounds gives, for \(R_T=A_0+\sqrt{2\log(4\lceil T\rceil)}\), \[\mathbb P(\sup_{[0,T]}|g|\le R_T)\ge\tfrac12.\] Thus \[\mathbb P(\inf_{[0,T]}(Z+g)>\ell)\ge\tfrac12 Q(\ell+R_T),
 \qquad -\log Q(\ell+R_T)=O(\log T).\] It follows directly that \(\theta_f^\ell=0\) at every fixed level. Its spectral measure has both a zero atom and a nonzero absolutely continuous component, satisfies \(\mathcal L\), and has infinite origin density. This agrees with Proposition 2.25 and Theorem 1, rather than refuting either.\par
These are peripheral range errors, with no demonstrated use of the overstated assertions in the core proofs. Accordingly they are Minor. Restricting Remark 3 to finite positive origin density avoids these counterexamples; its positive-level claim must still be presented as a conjecture, not as a proved theorem.\par
\section*{Coverage, counterevidence and limits}
The complete source was read in order, including all definitions, quantifiers, boundary cases, proofs and the bibliography. The singular-measure argument, the order of smooth approximation away from and near the origin, the change-of-level comparisons, and the compact-support condition in persistence sampling were followed. The seven problems above are the demonstrated issues retained after considering alternative readings. They are not a severity quota.\par
Several apparent problems do not warrant additional substantive findings. In Proposition 2.6 (p. 11), the final bound \((1-2/x)^x\ge e^{-2}\) is false, but the preceding, sharper Gaussian-tail bound retains \(2e^{-1/2}/x\) and suffices for the stated constant \(0.05\). The supremum direction in the displayed Slepian formulation (p. 13) is inconsistent with its subsequent use, whereas the operative infimum comparison has the correct direction. Claim 4.3 (p. 28) uses nonnested interval partitions as though they formed a filtration; differentiation of singular measures supplies the needed almost-everywhere conclusion. These observations do not supply independent counterexamples to those results.\par
The convexity argument on p. 33 can use a left secant at \(\ell/2\) to obtain the finite upper bound on the relevant derivative. The block count in Lemma 5.7 (pp. 36–38) must retain the factor \(1+\delta\); the resulting error is of the same order as the intended bound. Neither observation is presented as an additional proved failure of a final theorem. The broader statement of Proposition 5.1 includes smooth densities vanishing at zero, whereas its invocation of Lemma 5.7 uses the positive-origin-density hypothesis. The smooth approximants needed in the principal finite-positive-density argument do have that property, after normalizing the variance. This review does not independently establish all zero-density, negative-level cases of that broader proposition.\par
For Theorem 5, the distinction between ball sampling for \(\rho\in\mathcal L\) and persistence sampling for compactly supported \(\rho\in\mathcal L\cap\mathcal M\) was retained. The aliasing calculation in Proposition 7.2 (p. 47) was checked: the narrow spectral peaks have summable total mass, while their folded density at zero diverges by a harmonic sum. It is consistent with the importance of the compact-support condition and does not supply a new defect.\par
\subsection*{External mathematical inputs}
Royen’s primary preprint [royen], all seven pages, was read and visually checked for the symmetric-convex Gaussian correlation input used in the local repair of P02. This is a check of the applicable theorem, its degeneracy extension, and its role here; it is not an independent reproof of every multivariate-gamma identity in that paper.\par
For the spectral persistence estimates cited as source reference [28], the acquired primary text was [ffn], arXiv version 2 of 4 August 2020. Pages 1–5 and 8–25 were read and visually checked. In particular, its normalization and positive-moment convention, the absolutely continuous independent component, the integrated variance estimates, Theorems 4.1 and 5.1, Remark 2, and the proof of Proposition 3 were checked. This preprint has 31 pages; the 2021 journal article has a different pagination and was not acquired for line comparison. No equality of the two texts is asserted.\par
The observed Remark 2 in [ffn], p. 19, and its proof on p. 23, require the near-origin upper bound at all scales. This matters for a proposed replacement of the nonexistence example in P06: replacing the oscillation alone does not certify the subsequence persistence bounds, their normalization, or the precise applicability of the cited journal version. Those are explicit remaining repair tasks, not an additional version-independent finding established here. The main counterexamples in P01–P07 use direct analytic arguments, rather than an assumption of equivalence between reference versions.\par
\section*{Independent checks and disposition}
The reviewer-written calculation script evaluates normal-tail formulas, the independent-grid concentration lower bound, the removed spectral mass lower bound, the finite atomic covariance contradiction, and the oscillatory-integral expansion. It completed with exit code zero. No author code was executed. The calculations support the displayed analytic proofs; ordinary floating-point evaluations are not described as interval-certified numerical proofs. In particular, the stored remainder for the oscillatory integral bounds the exact analytic truncation error, not floating-point rounding.\par
The reviewed source has SHA256\par
> `289d568bfcfe53e56d54dcb01c2536d9c2c718aba211cfedd8cfcc698755cb04`\par
and its frozen source manifest has SHA256\par
> `d9ff41d9d197506de83917fa8c590d9c6e0cefc13764499f9b3e50448a3c9dbf`.\par
The exact TeX, compiled PDF, compiler, dependency payloads, logs, independent calculations, source-page images and actual inspection records are bound by the accompanying provenance and inventory. All scientific reading and judgments were made by the primary executor in this conversation. The helper performed only mechanical acquisition, binding, extraction and rendering.\par
The appropriate initial disposition is issues found: six Major and one Minor, with the scientific decision still pending. P01–P05 invalidate specified auxiliary statements or proof steps used in the principal arguments. P06 invalidates the advertised nonexistence derivation for its actual displayed density. P07 corrects peripheral ranges. None is presented as a demonstrated counterexample to Theorems 1–5. Local repairs proved here should be incorporated explicitly; repairs identified only as programs need a new mathematical check before acceptance.\par
The controller’s reported observation of an Academic/Pro user interface is retained as an attributed UI observation. Neither item is an independent backend certification or a proof of equivalence. Initial report completion does not constitute independent acceptance, and there are no human reviews recorded.\par
9 N. D. Feldheim, O. N. Feldheim and S. Mukherjee, Persistence and ball exponents for Gaussian stationary processes, arXiv:2112.04820v3 (3 April 2025), 51 pp. https://arxiv.org/abs/2112.04820v3. Associated journal DOI: https://doi.org/10.1002/cpa.22255. Wiley, article record for Persistence and Ball Exponents for Gaussian Stationary Processes, Commun. Pure Appl. Math. 78 (2025), 1949–2000. First online 28 April 2025; abstract correction noted 26 June 2025. https://onlinelibrary.wiley.com/doi/10.1002/cpa.22255. Metadata and correction notice checked 29 October 2026. T. Royen, A simple proof of the Gaussian correlation conjecture extended to some multivariate gamma distributions, arXiv:1408.1028v2 (13 August 2014), 7 pp. https://arxiv.org/abs/1408.1028v2. N. Feldheim, O. Feldheim and S. Nitzan, Persistence of Gaussian stationary processes: a spectral perspective, arXiv:1709.00204v2 (4 August 2020), 31 pp. https://arxiv.org/abs/1709.00204v2. Associated journal article: Ann. Probab. 49 (2021), 1067–1096, https://doi.org/10.1214/20-AOP1470.\par
\end{document}
